\section{AIME I 2008}
        \begin{enumerate}
        	\item (\textit{AIME I 2008}) Of the students attending a school party, $60\%$ of the students are girls, and $40\%$ of the students like to dance. After these students are joined by $20$ more boy students, all of whom like to dance, the party is now $58\%$ girls. How many students now at the party like to dance?


 
 \textbf{Answer: }252



\textbf{Solution 1a}

Say that there were $3k$ girls and $2k$ boys at the party originally. $2k$ like to dance. Then, there are $3k$ girls and $2k + 20$ boys, and $2k + 20$ like to dance.

 Thus, $\dfrac{3k}{5k + 20} = \dfrac{29}{50}$, solving gives $k = 116$. Thus, the number of people that like to dance is $2k + 20 = \boxed{252}$.

 



\textbf{Solution 1b}

Let $T$ be the original total. Then, there are $\frac{3T}{5}$ girls. Then, $\frac{\frac{3T}{5}}{T+20}=\frac{29}{50}$. Solving, $T=580$.

 $0.4\cdot580+20=\boxed{252}$.

 



\textbf{Solution 2}

Let the number of girls be $g$. Let the number of total people originally be $t$.

 We know that $\frac{g}{t}=\frac{3}{5}$ from the problem.

 We also know that $\frac{g}{t+20}=\frac{29}{50}$ from the problem.

 We now have a system and we can solve.

 The first equation becomes:

 $3t=5g$.

 The second equation becomes: 

 $50g=29t+580$

 Now we can sub in $30t=50g$ by multiplying the first equation by $10$. We can plug this into our second equation.

 $30t=29t+580$
$t=580$

 We know that there were originally $580$ people. Of those, $\frac{2}{5}*580=232$ like to dance.

 We also know that with these people, $20$ boys joined, all of whom like to dance. We just simply need to add $20$ to get $232+20=\boxed{252}$




\textbf{Solution 3}

Let $p$ denote the total number of people at the party. Then, because we know the proportions of boys to $p$ both before and after 20 boys arrived, we can create the following equation:
 \[0.4p+20 = 0.42(p+20)\]Solving for p gives us $p=580$, so the solution is $0.4p+20 = \boxed{252}$




\textbf{Solution 4 (Cheese)}

Assume all the boys like to dance and none of the girls like to dance. We then proceed like the previous solutions.

 


	\item (\textit{AIME I 2008}) Square $AIME$ has sides of length $10$ units.  Isosceles triangle $GEM$ has base $EM$, and the area common to triangle $GEM$ and square $AIME$ is $80$ square units.  Find the length of the altitude to $EM$ in $\triangle GEM$.


 
 \textbf{Answer: }025



\textbf{Solution 1}

Note that if the altitude of the triangle is at most $10$, then the maximum area of the intersection of the triangle and the square is $5\cdot10=50$.This implies that vertex G must be located outside of square $AIME$.

 
                    \begin{figure}[H]
                        \centering
                        \includegraphics[width=0.5\textwidth]{Solution/AIME_I_2008_2_Solution_Figure_1.png}
                    \end{figure}
                    

 Let $GE$ meet $AI$ at $X$ and let $GM$ meet $AI$ at $Y$. Clearly, $XY=6$ since the area of \href{https://artofproblemsolving.com/wiki/index.php?title=Trapezoid}{trapezoid} $XYME$ is $80$. Also, $\triangle GXY \sim \triangle GEM$.

 Let the height of $GXY$ be $h$. By the similarity, $\dfrac{h}{6} = \dfrac{h + 10}{10}$, we get $h = 15$. Thus, the height of $GEM$ is $h + 10 = \boxed{025}$.

 


	\item (\textit{AIME I 2008}) Ed and Sue bike at equal and constant rates.  Similarly, they jog at equal and constant rates, and they swim at equal and constant rates.  Ed covers $74$ kilometers after biking for $2$ hours, jogging for $3$ hours, and swimming for $4$ hours, while Sue covers $91$ kilometers after jogging for $2$ hours, swimming for $3$ hours, and biking for $4$ hours.  Their biking, jogging, and swimming rates are all whole numbers of kilometers per hour.  Find the sum of the squares of Ed's biking, jogging, and swimming rates.


 
 \textbf{Answer: }314



\textbf{Solution 1}

Let the biking rate be $b$, swimming rate be $s$, jogging rate be $j$, all in km/h.

 We have $2b + 3j + 4s = 74,2j + 3s + 4b = 91$. Subtracting the second from twice the first gives $4j + 5s = 57$. Mod 4, we need $s\equiv1\pmod{4}$. Thus, $(j,s) = (13,1),(8,5),(3,9)$.

 $(13,1)$ and $(3,9)$ give non-integral $b$, but $(8,5)$ gives $b = 15$. Thus, our answer is $15^{2} + 8^{2} + 5^{2} = \boxed{314}$.

 \textit{Note}: You can really quickly take the first equation $2b + 3j + 4s = 74$ mod 2, which will give $j \equiv 0 \pmod{2}$. The only one of the orde

 


	\item (\textit{AIME I 2008}) There exist unique positive integers $x$ and $y$ that satisfy the equation $x^2 + 84x + 2008 = y^2$. Find $x + y$.


 
 \textbf{Answer: }080



\textbf{Solution 1}

\href{https://artofproblemsolving.com/wiki/index.php?title=Completing\_the\_square}{Completing the square}, $y^2 = x^2 + 84x + 2008 = (x+42)^2 + 244$. Thus $244 = y^2 - (x+42)^2 = (y - x - 42)(y + x + 42)$ by \href{https://artofproblemsolving.com/wiki/index.php?title=Difference\_of\_squares}{difference of squares}. 

 Since $244$ is even, one of the factors is even. A \href{https://artofproblemsolving.com/wiki/index.php?title=Parity}{parity} check shows that if one of them is even, then both must be even. Since $244 = 2^2 \cdot 61$, the factors must be $2$ and $122$. Since $x,y > 0$, we have $y - x - 42 = 2$ and $y + x + 42 = 122$; the latter equation implies that $x + y = \boxed{080}$. 

 Indeed, by solving, we find $(x,y) = (18,62)$ is the unique solution. 

 



\textbf{Solution 2}

We complete the square like in the first solution: $y^2 = (x+42)^2 + 244$. Since consecutive squares differ by the consecutive odd numbers, we note that $y$ and $x+42$ must differ by an even number. We can use casework with the even numbers, starting with $y-(x+42)=2$.
 \begin{align*}2(x+42)+1+2(x+42)+3&=244\\ \Leftrightarrow x&=18\end{align*}


 Thus, $y=62$ and the answer is $\boxed{080}$.

 



\textbf{Solution 3}

We see that $y^2 \equiv x^2 + 4 \pmod{6}$. By \href{https://artofproblemsolving.com/wiki/index.php?title=Quadratic\_residue}{quadratic residues}, we find that either $x \equiv 0, 3 \pmod{6}$. Also, $y^2 \equiv (x+42)^2 + 244 \equiv (x+2)^2 \pmod{4}$, so $x \equiv 0, 2 \mod{4}$. Combining, we see that $x \equiv 0 \mod{6}$. 

 Testing $x = 6$ and other multiples of $6$, we quickly find that $x = 18, y = 62$ is the solution. $18+62=\boxed{080}$




\textbf{Solution 4}

We solve for x: $x^2 + 84x + 2008-y^2 = 0$
$x=\dfrac{-84+\sqrt{84^2-4\cdot 2008+4y^2}}{2}=-42+\sqrt{y^2-244}$

 So $y^2-244$ is a perfect square. Since 244 is even, the difference $\sqrt{y^2-244} -y^2$ is even, so we try $y^2-244=(y-2)^2$: $-244=-4y+4$, $y=62$.

 Plugging into our equation, we find that $x=18$, and $(x,y)=(18,62)$ indeed satisfies the original equation. $x+y=\boxed{080}$




\textbf{Solution 5}

Let $y=x+d$ for some $d>0$, substitute into the original equation to get $84x + 2008 = 2xd + d^2$.

 All terms except for the last one are even, hence $d^2$ must be even, hence let $d=2e$. We obtain $21x + 502 = xe + e^2$. Rearrange to $502-e^2 = x(e-21)$.

 Obviously for $0<e<21$ the right hand side is negative and the left hand side is positive. Hence $e\geq 21$. Let $e=21+f$, then $f\geq 0$.

 We have $502 - (21+f)^2 = xf$. Left hand side simplifies to $61 - 42f + f^2$. As $x$ must be an integer, $f$ must divide the left hand side. But $61$ is a prime, which only leaves two options: $f=1$ and $f=61$.

 Option $f=61$ gives us a negative $x$. Option $f=1$ gives us $x=61/f - 42 + f = 18$, and $y = x + d= x + 2e = x + 2(21+f) = 18 + 44 = 62$, hence $x+y=\boxed{080}$.

 


 



\textbf{Solution 6}

First complete the square to get $y^2 = (x+42)^2 + 244$. Remember that squares are the sums of consecutive odd integers, so when the difference between the two squares is 244, the two squares must be an even number of odd integers apart. However, there is only one distinct solution, as the problem states, and very quickly you will realize that only two odd integers work. When there are four, then the numbers are not odd, and when it is any other even integer it does not divide. So we need two consecutive odd integers that sum to 244. Easily we find 121 and 123. 121 is the 61st odd integer and 123 is the 62nd odd integer, so $(x+42)^2$ is the sum of the first 60 odd integers, or $(60)^2$, while $y^2$ is $62^2$ for the same reasons. That way we get $x=12$, $y=62$, hence $x+y=\boxed{080}$.

 
-jackshi2006

 


	\item (\textit{AIME I 2008}) A right circular cone has base radius $r$ and height $h$.  The cone lies on its side on a flat table.  As the cone rolls on the surface of the table without slipping, the point where the cone's base meets the table traces a circular arc centered  at the point where the vertex touches the table.  The cone first returns to its original position on the table after making $17$ complete rotations.  The value of $h/r$ can be written in the form $m\sqrt {n}$, where $m$ and $n$ are positive integers and $n$ is not divisible by the square of any prime.  Find $m + n$.


 
 \textbf{Answer: }014



\textbf{Solution}

The path is a circle with radius equal to the slant height of the cone, which is $\sqrt {r^{2} + h^{2}}$. Thus, the length of the path is $2\pi\sqrt {r^{2} + h^{2}}$.

 Also, the length of the path is 17 times the circumference of the base, which is $34r\pi$. Setting these equal gives $\sqrt {r^{2} + h^{2}} = 17r$, or $h^{2} = 288r^{2}$. Thus, $\dfrac{h^{2}}{r^{2}} = 288$, and $\dfrac{h}{r} = 12\sqrt {2}$, giving an answer of $12 + 2 = \boxed{014}$.

 


	\item (\textit{AIME I 2008}) A triangular array of numbers has a first row consisting of the odd integers $1,3,5,\ldots,99$ in increasing order.  Each row below the first has one fewer entry than the row above it, and the bottom row has a single entry.  Each entry in any row after the top row equals the sum of the two entries diagonally above it in the row immediately above it.  How many entries in the array are multiples of $67$?


 
 \textbf{Answer: }017



\textbf{Solution 1}

Let the $k$th number in the $n$th row be $a(n,k)$. Writing out some numbers, we find that $a(n,k) = 2^{n-1}(n+2k-2)$.<sup>[\href{https://artofproblemsolving.com\#note\_1}{1}]</sup>

 We wish to find all $(n,k)$ such that $67| a(n,k) = 2^{n-1} (n+2k-2)$. Since $2^{n-1}$ and $67$ are \href{https://artofproblemsolving.com/wiki/index.php?title=Relatively\_prime}{relatively prime}, it follows that $67|n+2k-2$. Since every row has one less element than the previous row, $1 \le k \le 51-n$ (the first row has $50$ elements, the second $49$, and so forth; so $k$ can range from $1$ to $50$ in the first row, and so forth). Hence 

 $n+2k-2 \le n + 2(51-n) - 2 = 100 - n \le 100,$

 it follows that $67| n - 2k + 2$ implies that $n-2k+2 = 67$ itself.

 Now, note that we need $n$ to be odd, and also that $n+2k-2 = 67 \le 100-n \Longrightarrow n \le 33$.

 We can check that all rows with odd $n$ satisfying $1 \le n \le 33$ indeed contains one entry that is a multiple of $67$, and so the answer is $\frac{33+1}{2} = \boxed{017}$.

 



 \href{https://artofproblemsolving.com\#ref\_1}{\textasciicircum{}} Proof: Indeed, note that $a(1,k) = 2^{1-1}(1+2k-2)=2k-1$, which is the correct formula for the first row. We claim the result by \href{https://artofproblemsolving.com/wiki/index.php?title=Induction}{induction} on $n$. By definition of the array, $a(n,k) = a(n-1,k)+a(n-1,k+1)$, and by our inductive hypothesis, 
 \begin{align*}a(n,k) &= a(n-1,k)+a(n-1,k+1)\\ &= 2^{n-2}(n-1+2k-2)+2^{n-2}(n-1+2(k+1)-2)\\&=2^{n-2}(2n+4k-4)\\&=2^{n-1}(n+2k-2)\end{align*}
thereby completing our induction.

 



\textbf{Solution 2}

The result above is fairly intuitive if we write out several rows and then divide all numbers in row $r$ by $2^{r-1}$ (we can do this because dividing by a power of 2 doesn't affect divisibility by $67$). The second row will be $2, 4, 6, \cdots , 98$, the third row will be $3, 5, \cdots, 97$, and so forth. Clearly, only the odd-numbered rows can have a term divisible by $67$. However, with each row the row will have one less element, and the $99-67+1 = 33$rd row is the last time $67$ will appear. Therefore the number of multiples of 67 in the entire array is $\frac{33+1}{2} = \boxed{017}$.

 



\textbf{Solution 3}


 \[\begin{array}{ccccccccccccccc} \color{blue}1 & & \color{blue}3 & & \color{blue}5 & & \color{blue}7 & & \color{blue}9 & & \color{blue}11 & & \color{blue}13 & & \color{blue}\dots \\ & \color{red}4 & & \color{red}8 & & \color{red}12 & & \color{red}16 & & \color{red}20 & & \color{red}24 & & \color{red}\dots & \\ & & \color{blue}12 & & \color{blue}20 & & \color{blue}28 & & \color{blue}36 & & \color{blue}44 & & \color{blue}\dots & & \\ & & & \color{red}32 & & \color{red}48 & & \color{red}64 & & \color{red}80 & & \color{red}\dots & & & \\ & & & & \color{blue}80 & & \color{blue}112 & & \color{blue}144 & & \color{blue}\dots & & & & \\ & & & & & \color{red}192 & & \color{red}256 & & \color{red}\dots & & & & &  \end{array}\]

 We begin by listing a smaller portion of the given triangular array, as shown in the diagram. Observe that for any term \(k\) in the \(n\)th row, the corresponding term in the \((n+2)\)th row is \(4k\), and that term is also one position closer to the edge of the triangle. Since \(67\) is prime, the only way a term in the array can be a multiple of \(67\) is if it originates from the \(67\) in the first row. Next, note that there cannot be any multiples of \(67\) in the even-numbered rows (rows \(2,4,6,\dots,50\)). In particular, the second row consists entirely of multiples of \(4\) up to \(196\), none of which contain a factor of \(67\). Since each even row is obtained by repeatedly multiplying by \(4\), no even row can contain a multiple of \(67\). Now observe that \(67\) is the \(17\)th term from the right edge of the triangle. Each time a term moves two rows down, it shifts one position closer to the edge. Once a term reaches the edge, it is no longer multiplied by \(4\), so it cannot produce any new multiples of \(67\). Therefore, there are exactly \(\boxed{017}\) multiples of \(67\) in the array.

 



\textbf{Solution 4 (very similar to 1, different beginning and ending)}

Like in solution 3, experiment with a couple of the terms. Paying attention to the first couple terms, the pattern for EVERY term can soon come to life. Denote the rows, top to bottom, as $t = 1, 2, \dots ,50$. Denote $n$ to be the $n$th term of the sequence. Then, every term satisfies

 
 \[t(2^{t-1}) + (n-1)2^{t} \ : \ 1 \le t \le 50, 1 \le n \le 51-t\].

 This can easily be proved via induction, as shown in solution 1.

 From here, factoring out shows $2^{t-1} (t + 2n - 2) \pmod{67} \equiv 0$. Because 67 is prime, it follows that 67 divides $t+2n-2$. We have an upper bound for this however, particularly when $n = 51 - t$, in which we have $100-t$ from $t+2n-2$, and because $t \ge 1$, minimal value obtained is 99. Because of that, either it is equivalent to 0 or 67. If it is equivalent to 0, then we have $2-t = 2n$ and so $1-\frac{t}{2} = n$. For $n$ to be an integer in the range, the only $t$ that works is $t = 2$. However, $n$ would equal 0, which is not allowed. So no cases here. Otherwise, we have that the value itself is 67, in which we have $69-t = 2n$, so $t$ must be odd. Through definition, we have $\frac{69-t}{2} \le 51-t$, solving yields $69-t \le 102 - 2t$, so $t \le 33$. There are 17 such odd numbers.

 Because no other value $t$ will satisfy those prerequisites, only $\boxed{017}$ multiples are viable.

 


	\item (\textit{AIME I 2008}) Let $S_i$ be the set of all integers $n$ such that $100i\leq n < 100(i + 1)$.  For example, $S_4$ is the set ${400,401,402,\ldots,499}$.  How many of the sets $S_0, S_1, S_2, \ldots, S_{999}$ do not contain a perfect square?


 
 \textbf{Answer: }708



\textbf{Solution}

The difference between consecutive squares is $(x + 1)^2 - x^2 = 2x + 1$, which means that all squares above $50^2 = 2500$ are more than $100$ apart.

 Then the first $26$ sets ($S_0,\cdots S_{25}$) each have at least one perfect square because the differences between consecutive squares in them are all less than $100$. Also, since $316$ is the largest $x$ such that $x^2 < 100000$ ($100000$ is the upper bound which all numbers in $S_{999}$ must be less than), there are $316 - 50 = 266$ other sets after $S_{25}$ that have a perfect square.

 There are $1000 - 266 - 26 = \boxed{708}$ sets without a perfect square.

 


	\item (\textit{AIME I 2008}) Find the positive integer $n$ such that 


 
 \[\arctan\frac {1}{3} + \arctan\frac {1}{4} + \arctan\frac {1}{5} + \arctan\frac {1}{n} = \frac {\pi}{4}.\]

 \textbf{Answer: }047



\textbf{Solution 1}

Since we are dealing with acute angles, $\tan(\arctan{a}) = a$.

 Note that $\tan(\arctan{a} + \arctan{b}) = \dfrac{a + b}{1 - ab}$, by tangent addition. Thus, $\arctan{a} + \arctan{b} = \arctan{\dfrac{a + b}{1 - ab}}$.

 Applying this to the first two terms, we get $\arctan{\dfrac{1}{3}} + \arctan{\dfrac{1}{4}} = \arctan{\dfrac{7}{11}}$.

 Now, $\arctan{\dfrac{7}{11}} + \arctan{\dfrac{1}{5}} = \arctan{\dfrac{23}{24}}$.

 We now have $\arctan{\dfrac{23}{24}} + \arctan{\dfrac{1}{n}} = \dfrac{\pi}{4} = \arctan{1}$. Thus, $\dfrac{\dfrac{23}{24} + \dfrac{1}{n}}{1 - \dfrac{23}{24n}} = 1$; and simplifying, $23n + 24 = 24n - 23 \Longrightarrow n = \boxed{47}$.

 



\textbf{Solution 2 (generalization)}

From the expansion of $e^{iA}e^{iB}e^{iC}e^{iD}$, we can see that
 \[\cos(A + B + C + D) = \cos A \cos B \cos C \cos D - \frac{1}{4} \sum_{\textrm sym} \sin A \sin B \cos C \cos D + \sin A \sin B \sin C \sin D,\]and
 \[\sin(A + B + C + D) = \sum_{\textrm cyc}\sin A \cos B \cos C \cos D - \sum_{\textrm cyc} \sin A \sin B \sin C \cos D .\]If we divide both of these by $\cos A \cos B \cos C \cos D$, then we have
 \[\tan(A + B + C + D) = \frac {1 - \sum \tan A \tan B + \tan A \tan B \tan C \tan D}{\sum \tan A - \sum \tan A \tan B \tan C},\]which makes for more direct, less error-prone computations. Substitution gives the desired answer.

 



\textbf{Solution 3: Complex Numbers}

Adding a series of angles is the same as multiplying the complex numbers whose arguments they are. In general, $\arctan\frac{1}{n}$, is the argument of $n+i$. The sum of these angles is then just the argument of the product 
 \[(3+i)(4+i)(5+i)(n+i)\]and expansion give us $(48n-46)+(48+46n)i$. Since the argument of this complex number is $\frac{\pi}{4}$, its real and imaginary parts must be equal; then, we can we set them equal to get
 \[48n - 46 = 48 + 46n.\]Therefore, $n=\boxed{47}$.

 



\textbf{Solution 4 Sketch}

You could always just bash out $\sin(a+b+c)$ (where a,b,c,n are the angles of the triangles respectively) using the sum identities again and again until you get a pretty ugly radical and use a triangle to get $\cos(a+b+c)$ and from there you use a sum identity again to get $\sin(a+b+c+n)$ and using what we found earlier you can find $\tan(n)$ by division that gets us $\frac{23}{24}$




\textbf{Solution 5 (Deriving formulae)}

Define $\operatorname{arctan}(x) = A$ and $\operatorname{arctan}(y) = B$. We will attempt to find $\arctan(x) + \arctan(y)$. That is, we attempt to find $A+B$. Notice $\arctan(x) = A$ implies $\tan(A) = x$, and similarly for $\tan(B) = y$. We know $\tan(A+B) = \frac{\tan(A) + \tan(B)}{1-\tan(A)\tan(B)}$, so $\tan(A+B) = \frac{x+y}{1-xy}$. In addition, we know that arc tangent both sides yields $A+B = \arctan(\frac{x+y}{1-xy})$. Therefore, $\arctan(x) + \arctan(y) = \arctan(\frac{x+y}{1-xy})$. Applying that with $x = \frac{1}{3}$ and $y = \frac{1}{4}$, we have $\arctan(\frac{7}{11})$. Applying this with $\arctan(\frac{1}{5})$, we have $\arctan(\frac{23}{24})$. Finally, applying that with $\arctan(\frac{1}{n})$, we have $\arctan(\frac{23n+24}{24n-23}) = \frac{\pi}{4}$. Notice that $\arctan(1) = \frac{\pi}{4}$, so we have $23n + 24 = 24n-23$, or $n = \boxed{047}$.

 


	\item (\textit{AIME I 2008}) Ten identical crates each of dimensions $3$ ft $\times$ $4$ ft $\times$ $6$ ft.  The first crate is placed flat on the floor.  Each of the remaining nine crates is placed, in turn, flat on top of the previous crate, and the orientation of each crate is chosen at random.  Let $\frac {m}{n}$ be the probability that the stack of crates is exactly $41$ ft tall, where $m$ and $n$ are relatively prime positive integers.  Find $m$.


 
 \textbf{Answer: }190



\textbf{Solution}

Only the heights matter, and each crate is either 3, 4, or 6 feet tall with equal probability. We have the following:

 
 \begin{align*}3a + 4b + 6c &= 41\\ a + b + c &= 10\end{align*}


 Subtracting 3 times the second from the first gives $b + 3c = 11$, or $(b,c) = (2,3),(5,2),(8,1),(11,0)$. The last doesn't work, obviously. This gives the three solutions $(a,b,c) = (5,2,3),(3,5,2),(1,8,1)$. In terms of choosing which goes where, the first two solutions are analogous.

 For $(5,2,3),(3,5,2)$, we see that there are $2\cdot\dfrac{10!}{5!2!3!} = 10\cdot9\cdot8\cdot7$ ways to stack the crates. For $(1,8,1)$, there are $\dfrac{10!}{8!1!1!} = 90$. Also, there are $3^{10}$ total ways to stack the crates to any height.

 Thus, our probability is $\dfrac{10\cdot9\cdot8\cdot7 + 90}{3^{10}} = \dfrac{10\cdot8\cdot7 + 10}{3^{8}} = \dfrac{570}{3^8} = \dfrac{190}{3^{7}}$. Our answer is the numerator, $\boxed{190}$.

 \textbf{1 Min Solution} It would be helpful for this solution to be reformatted. To start with, let us observe the three numbers. Note that $3$ and $6$ are both divisible by $3$, so the number of $4$-crates must be congruent to $41\bmod{3}$, which is also congruent to $2\bmod{3}$. Our solutions for the number of $4$-crates will repeat mod $3$, so if $x$ is a solution, so is $x+3$. By inspection, we have that $2$ is solution, and so are $5$ and $8$. Each solution splits into its own case.We must solve the equation $41-4z=6x+3y$, simultaneously with $x+y=10-z$. Note that we already know the possible values of $z$. Solving these (it's AIME $9$, you should be able to do this and if anyone feels like they want to write a rundown of this please go ahead), we get the solution  sets $\{8,1,1\},\{5,2,3\},$ and $\{2,3,5\}$. We can count the number of possible arrangements for each solution by taking $\dbinom{10}{z}$ and then multiplying by $\dbinom{10-z}{x}$ (the solution sets, for the sake of consistency, are in the form $z,x,y$). Summing the results for all the solutions gives us $5130$. Finally, to calculate the probability we must determine our denominator. Since we have $3$ ways to arrange each block, our denominator is $3^{10}$. $\frac{5130}{3^{10}}=\frac{190}{3^7}$. The answer is $m=\boxed{190}$.

 



\textbf{Solution 2 (a more direct approach)}

Let's make two observations. We are trying to find the number of ways we can add $3\text{s}, 4\text{s}$, and $6\text{s}$ to get $41$, and the total number of (non-distinct) sums possible is $3^{10}$.Then we just use casework to easily and directly solve for the number of ways to get $41$. To begin, the minimum sum is produced with $10$ threes, so WLOG we can solve for the number of ways to get $11$ with $0\text{s}, 1\text{s}$, and $3\text{s}$. 

 Case I: $0$ zeroes, $0$ threes, $11$ onesImpossible, because there are only ten available spots.

 Case II: $1$ zero, $1$ three, $8$ onesThis is just $\frac{10!}{8!}$, so there are $90$ possibilities.

 Case III: $3$ zeroes, $2$ threes, $5$ onesThis is just $\frac{10!}{3!2!5!}$. This gives $2520$ possibilities.

 Case IV: 5 zeroes, 3 threes, and 2 ones.This is the same as case $3$, so also $2520$ possibilities.

 $90+2520+2520=5130$
$5130$ has three powers of $3$, so $5130$ divided by $27$ is $\boxed{190}$.

 -jackshi2006

 



\textbf{Solution 3 (Generating Functions)}

Note we are placing 10 crates where each "height" is 3, 4, 6 and we want all the heights to sum to 41. We can model this as the generating function 
 \[\left(x^3+x^4+x^6\right)^{10}\] where we want the coefficient of $x^{41}.$First off, factor this to get 
 \[{x^{30} \left(1 + x + x^3\right)^{10}}\] and then see that we want the coefficient of $x^{11}$ in $\left(1+x+x^3\right)^{10}.$From multinomial theorem, this expansion is 
 \[\sum_{a+b+c=10}\binom{10}{a,b,c}1^ax^bx^{3c}\]If we want the coefficient of $x^{11}$ then we need $b + 3c = 11.$ with $b + c \le 10$(from the multinomial expansion).This has the solutions $(b, c) = \{8, 1\}, \{5, 2\}, \{2, 3\}.$ Note that the denominator of the answer is just $3^{10}$ since there are 3 ways to orientate every crate and there are 10 creates. Thus, our answer is 
 \[\frac{\binom{10}{8,1,1} + \binom{10}{5,2,3} + \binom{10}{2,3,5}}{3^{10}} = \frac{190}{3^7} \rightarrow \boxed{190}\]



	\item (\textit{AIME I 2008}) Let $ABCD$ be an isosceles trapezoid with $\overline{AD}||\overline{BC}$ whose angle at the longer base $\overline{AD}$ is $\dfrac{\pi}{3}$. The diagonals have length $10\sqrt {21}$, and point $E$ is at distances $10\sqrt {7}$ and $30\sqrt {7}$ from vertices $A$ and $D$, respectively. Let $F$ be the foot of the altitude from $C$ to $\overline{AD}$. The distance $EF$ can be expressed in the form $m\sqrt {n}$, where $m$ and $n$ are positive integers and $n$ is not divisible by the square of any prime. Find $m + n$.


 
 \textbf{Answer: }032



\textbf{Solution 1}





	\item (\textit{AIME I 2008}) Consider sequences that consist entirely of $A$'s and $B$'s and that have the property that every run of consecutive $A$'s has even length, and every run of consecutive $B$'s has odd length. Examples of such sequences are $AA$, $B$, and $AABAA$, while $BBAB$ is not such a sequence. How many such sequences have length 14?


 
 \textbf{Answer: }172



\textbf{Solution 1a}

Let $a_n$ and $b_n$ denote, respectively, the number of sequences of length $n$ ending in $A$ and $B$. If a sequence ends in an $A$, then it must have been formed by appending two $A$s to the end of a string of length $n-2$. If a sequence ends in a $B,$ it must have either been formed by appending one $B$ to a string of length $n-1$ ending in an $A$, or by appending two $B$s to a string of length $n-2$ ending in a $B$. Thus, we have the \href{https://artofproblemsolving.com/wiki/index.php?title=Recursion}{recursions}
 \begin{align*} a_n &= a_{n-2} + b_{n-2}\\ b_n &= a_{n-1} + b_{n-2}  \end{align*}
By counting, we find that $a_1 = 0, b_1 = 1, a_2 = 1, b_2 = 0$. 
 \[\begin{array}{|r||r|r|||r||r|r|} \hline n & a_n & b_n & n & a_n & b_n\\ \hline 1&0&1& 8&6&10\\ 2&1&0& 9&11&11\\ 3&1&2& 10&16&21\\ 4&1&1& 11&22&27\\ 5&3&3& 12&37&43\\ 6&2&4& 13&49&64\\ 7&6&5& 14&80&92\\ \hline \end{array}\]Therefore, the number of such strings of length $14$ is $a_{14} + b_{14} = \boxed{172}$.

 



\textbf{Solution 1b}

Let $a_n$ and $b_n$ denote, respectively, the number of sequences of length $n$ ending in $A$ and $B$. 

 Additionally, let $t_n$ denote the total number of sequences of length $n$. Then, $t_n=a_n+b_n$, as the total amount of sequences of length $n$ consists of the sequences of length $n$ ending in $A$ and the sequences of length $n$ ending in $B$.
 \begin{align*} a_n &= a_{n-2} + b_{n-2}\\ b_n &= a_{n-1} + b_{n-2}  \end{align*}
The recursion for $a_n$ tells us that $a_n=a_{n-2}+b_{n-2}$. However, this is also the definition for $t_{n-2}$. Therefore, $a_n=t_{n-2}$.

 We also know from our recursion for $b_n$ that $b_n=a_{n-1}+b_{n-2}$. Substituting for $a_n$ and $b_n$ into our recursion for $t_n$ gives us $t_n=t_{n-2}+a_{n-1}+b_{n-2}$.

 Furthermore, note that since $a_n=t_{n-2}$, $a_{n-1}=t_{n-3}$. Furthermore, using our definition for $t_{n-2}$, we can rewrite $b_{n-2}$ as $t_{n-2}-a_{n-2}$. Substituting for $a_{n-1}$ and $b_{n-2}$ into our recursion for $t_n$ gives us $t_n=t_{n-2}+t_{n-3}+t_{n-2}-a_{n-2}$.

 Finally, note that since $a_n=t_{n-2}$, $a_{n-2}=t_{n-4}$. Substituting for $a_{n-2}$ into our recursion for $t_n$ gives us $t_n=2t_{n-2}+t_{n-3}-t_{n-4}$. We now have a recursion only in terms of $t$.

 By counting, we find that $t_1=1$, $t_2=1$, $t_3=3$, and $t_4=2$.
 \[\begin{array}{|r|r||r|r|}	 \hline	 n & t_n & n & t_n\\	 \hline	 1&1&8&16\\	 2&1&9&22\\ 3&3&10&37\\	 4&2&11&49\\	 5&6&12&80\\	 6&6&13&113\\	 7&11&14&172\\	 \hline \end{array}\]

 Therefore, the number of such sequences of length 14 is $\boxed{172}$.

 



\textbf{Solution 1c}

Define $A_{n}$ being sequences with the first letter being A, $B_{n}$ being sequences with the first letter being B. 

 
$A_{14}$ = $B_{12}$ + $B_{10}$ + $B_{8}$ + $B_{6}$ + $B_{4}$ + $B_{2}$ (1)$B_{14}$ = $A_{13}$ + $A_{11}$ + $A_{9}$ + $A_{7}$ + $A_{5}$ + $A_{3}$ + $A_{1}$ (2)

 
We also know that 

 
$A_{12}$ = $B_{10}$ + $B_{8}$ + $B_{6}$ + $B_{4}$ + $B_{2}$ $B_{12}$ = $A_{11}$ + $A_{9}$ + $A_{7}$ + $A_{5}$ + $A_{3}$ + $A_{1}$

 
We can Plug in $A_{12}$ and $B_{12}$ into (1) and (2) to get

 
$A_{14}$ = $B_{12}$ + $A_{12}$$B_{14}$ = $A_{13}$ + $B_{12}$

 
We have the equation and the rest of the solutions are in Solution 1a.

 Note: Only solution 1c is written by me. 1a and 1b are not written by me.

 



\textbf{Solution 2}

We replace "14" with "$2n$". We first note that we must have an even number of chunks of $B$'s, because of parity issues.  We then note that every chunk of $B$'s except the last one must end in the sequence $BAA$, since we need at least two $A$'s to separate it from the next chunk of $B$'s. The last chunk of $B$'s must, of course, end with a $B$.  Thus our sequence must look like this :
 \[\boxed{\quad A\text{'s} \quad} \boxed{\quad B\text{'s} \quad} BAA \boxed{\quad A\text{'s} \quad} \cdots \boxed{\quad B\text{'s} \quad}BAA \boxed{\quad A\text{'s} \quad} \boxed{\quad B\text{'s} \quad} B \boxed{\quad A\text{'s} \quad} ,\]where each box holds an even number of letters (possibly zero).

 If we want a sequence with $2k$ chunks of $B$'s, then we have $(2n - 6k + 2)$ letters with which to fill the boxes. Since each box must have an even number of letters, we may put the letters in the boxes in pairs. Then we have $(n - 3k + 1)$ pairs of letters to put into $4k + 1$ boxes. By a classic \href{https://artofproblemsolving.com/wiki/index.php?title=Balls-and-bins}{balls-and-bins} argument, the number of ways to do this is
 \[\binom{n + k + 1}{4k} .\]It follows that the total number of desirable sequences is
 \[\sum_k \binom{n + k + 1}{4k} .\]For $n = 7$, this evaluates as $\binom{8}{0} + \binom{9}{4} + \binom{10}{8} = 1 + 126 + 45 = \boxed{172}$.

 



\textbf{Solution 3}

There must be an even amount of runs of consecutive $B$s due to parity. Thus, we can split this sequence into the following cases: $A$, $BAAB$, $AABAAB$, $BAABAA$, $AABAABAA$, $BAABAABAAB$, $AABAABAABAAB$, $BAABAABAABAA$, and $AABAABAABAABAA$, in which the amount of letters in one run does not necessarily represent the amount of letters there can be.

 For the first case and the last case, there is only one possible sequence of letters.

 For all other cases, we can insert two of the same letter at a time into a run that has the exact same letter. For example, for the second case, we can insert two $A$s and make the sequence $BAAAAB$. There are three "slots" in which we can insert two additional letters in, and we must insert five groups of new letters. By \href{https://artofproblemsolving.com/wiki/index.php?title=Stars\_and\_bars}{stars and bars}, the number of ways for the second case is $\binom{7}{2}=21$.

 Applying this logic to all of the other cases gives us $\binom{7}{3}$, $\binom{7}{3}$, $\binom{7}{4}$, $\binom{8}{6}$, $\binom{8}{1}$, and $\binom{8}{1}$. Adding 1+$\binom{7}{2}$+$\binom{7}{3}$+$\binom{7}{3}$+$\binom{7}{4}$+$\binom{8}{6}$+$\binom{8}{1}$+$\binom{8}{1}$ gives us the answer $\boxed{172}$.

 


	\item (\textit{AIME I 2008}) On a long straight stretch of one-way single-lane highway, cars all travel at the same speed and all obey the safety rule: the distance from the back of the car ahead to the front of the car behind is exactly one car length for each 15 kilometers per hour of speed or fraction thereof (Thus the front of a car traveling 52 kilometers per hour will be four car lengths behind the back of the car in front of it.) A photoelectric eye by the side of the road counts the number of cars that pass in one hour. Assuming that each car is 4 meters long and that the cars can travel at any speed, let $M$ be the maximum whole number of cars that can pass the photoelectric eye in one hour. Find the quotient when $M$ is divided by 10.


 
 \textbf{Answer: }375



\textbf{Solution 1}

Let $n$ be the number of car lengths that separates each car (it is easy to see that this should be the same between each pair of consecutive cars.) Then their speed is at most $15n$. Let a \textit{unit} be the distance between the cars (front to front). Then the length of each unit is $4(n + 1)$. To maximize, in a unit, the CAR comes first, THEN the empty space. So at time zero, the car is right at the eye.

 Hence, we count the number of units that pass the eye in an hour: $\frac {15,000n\frac{\text{meters}}{\text{hour}}}{4(n + 1)\frac{\text{meters}}{\text{unit}}} = \frac {15,000n}{4(n + 1)}\frac{\text{units}}{\text{hour}}$. We wish to maximize this.

 Observe that as $n$ gets larger, the $+ 1$ gets less and less significant, so we take the limit as $n$ approaches infinity

 $\lim_{n\rightarrow \infty}\frac {15,000n}{4(n + 1)} = \lim_{n\rightarrow \infty}\frac {15,000}{4} = 3750$

 Now, as the speeds are clearly finite, we can never actually reach $3750$ full UNITs. However, we only need to find the number of CARS. We can increase their speed so that the camera stops (one hour goes by) after the car part of the $3750$th unit has passed, but not all of the space behind it. Hence, $3750$ cars is possible, and the answer is $\boxed {375}$.

 



\textbf{Solution 2}

Disclaimer: This is for the people who may not understand calculus, and is also how I did it.First, we assume several things. First, we assume the speeds must be multiples of 15 to maximize cars, because any less will be a waste. Second, we start with one car in front of the photoelectric eye.We first set the speed of the cars as $15k$. Then, the distance between them is $\frac{4}{1000} \times k\text{km}$. Therefore, it takes the car closest to the eye not on the eye $\frac{\frac{k}{250}}{15k}$ hours to get to the eye. There is one hour, so the amount of cars that can pass is $\frac{1}{\frac{\frac{k}{250}}{15k}}$, or $3750$ cars. When divided by ten, you get the quotient of $\boxed{375}$




\textbf{Solution 3}

Call the speed of each car $s$, so that the distance between each car will be $\lfloor a/15 \rfloor$. To maximize the number of cars that pass, let the first car back end line up with the eye. Notice that every $4/1000+ (4/1000)(\lfloor a/15 \rfloor)$ kilometers (distance from front of one car to the front of the consecutive car), one car will be captured. The time it takes to travel this distance is $\frac{4/1000+(4/1000)(\lfloor a/15 \rfloor)}{a}$, and the number of these intervals in an hour is $\frac{1}{\frac{4/1000+(4/1000)(\lfloor a/15 \rfloor)}{a}}=\frac{250a}{1+\lfloor a/15 \rfloor}$. Now, let $a=15n$ and we can proceed with Solution 1.

 


	\item (\textit{AIME I 2008}) Let


 $p(x,y) = a_0 + a_1x + a_2y + a_3x^2 + a_4xy + a_5y^2 + a_6x^3 + a_7x^2y + a_8xy^2 + a_9y^3$.


 Suppose that 


 $p(0,0) = p(1,0) = p( - 1,0) = p(0,1) = p(0, - 1) = p(1,1) = p(1, - 1) = p(2,2) = 0$.


 There is a point $\left(\frac {a}{c},\frac {b}{c}\right)$ for which $p\left(\frac {a}{c},\frac {b}{c}\right) = 0$ for all such polynomials, where $a$, $b$, and $c$ are positive integers, $a$ and $c$ are relatively prime, and $c > 1$.  Find $a + b + c$.


 
 \textbf{Answer: }040



\textbf{Solution 1}


 \begin{align*} p(0,0) &= a_0 \\ &= 0 \\ p(1,0) &= a_0 + a_1 + a_3 + a_6 \\ &= a_1 + a_3 + a_6 \\ &= 0 \\ p(-1,0) &= -a_1 + a_3 - a_6 \\ &= 0 \end{align*}


 Adding the above two equations gives $a_3 = 0$, and so we can deduce that $a_6 = -a_1$.

 Similarly, plugging in $(0,1)$ and $(0,-1)$ gives $a_5 = 0$ and $a_9 = -a_2$. Now,
 \begin{align*} p(1,1) &= a_0 + a_1 + a_2 + a_3 + a_4 + a_5 + a_6 + a_7 + a_8 + a_9 \\ &= 0 + a_1 + a_2 + 0 + a_4 + 0 - a_1 + a_7 + a_8 - a_2 \\ &= a_4 + a_7 + a_8 \\ &= 0 \\ p(1,-1) &= a_0 + a_1 - a_2 + 0 - a_4 + 0 - a_1 - a_7 + a_8 + a_2 \\ &= -a_4 - a_7 + a_8 \\ &= 0 \end{align*}
Therefore $a_8 = 0$ and $a_7 = -a_4$. Finally,
 \begin{align*} p(2,2) &= 0 + 2a_1 + 2a_2 + 0 + 4a_4 + 0 - 8a_1 - 8a_4 +0 - 8a_2 \\ &= -6 a_1 - 6 a_2 - 4 a_4 \\ &= 0 \end{align*}
So, $3a_1 + 3a_2 + 2a_4 = 0$, or equivalently $a_4 = -\frac{3(a_1 + a_2)}{2}$.

 Substituting these equations into the original polynomial $p$, we find that at $\left(\frac{a}{c}, \frac{b}{c}\right)$,
 \[a_1x + a_2y + a_4xy - a_1x^3 - a_4x^2y - a_2y^3 = 0 \iff\]
 \[a_1x + a_2y - \frac{3(a_1 + a_2)}{2}xy - a_1x^3 + \frac{3(a_1 + a_2)}{2}x^2y - a_2y^3 = 0 \iff\]
 \[a_1x(x - 1)\left(x + 1 - \frac{3}{2}y\right) + a_2y\left(y^2 - 1 - \frac{3}{2}x(x - 1)\right) = 0\].The remaining coefficients $a_1$ and $a_2$ are now completely arbitrary because the original equations impose no more restrictions on them. Hence, for the final equation to hold for all possible $p$, we must have $x(x - 1)\left(x + 1 - \frac{3}{2}y\right) = y\left(y^2 - 1 - \frac{3}{2}x(x - 1)\right) = 0$. 

 As the answer format implies that the $x$-coordinate of the root is non-integral, $x(x - 1)\left(x + 1 - \frac{3}{2}y\right) = 0 \iff x + 1 - \frac{3}{2}y = 0 \iff y = \frac{2}{3}(x + 1)\ (1)$. The format also implies that $y$ is positive, so $y\left(y^2 - 1 - \frac{3}{2}x(x - 1)\right) = 0 \iff y^2 - 1 - \frac{3}{2}x(x - 1) = 0\ (2)$. Substituting $(1)$ into $(2)$ and reducing to a quadratic yields $(19x - 5)(x - 2) = 0$, in which the only non-integral root is $x = \frac{5}{19}$, so $y = \frac{16}{19}$.

 The answer is $5 + 16 + 19 = \boxed{040}$.

 
                    \begin{figure}[H]
                        \centering
                        \includegraphics[width=0.5\textwidth]{Solution/AIME_I_2008_13_Solution_Figure_1.png}
                    \end{figure}
                    




\textbf{Solution 2}

Consider the cross section of $z = p(x, y)$ on the plane $z = 0$. We realize that we could construct the lines/curves in the cross section such that their equations multiply to match the form of $p(x, y)$ (same degree of $x$ and $y$ in terms) and they include the eight given points. One simple way to do this would be to use the equations $x = 0$, $x = 1$, and $y = \frac{2}{3}x + \frac{2}{3}$, giving us 

 $p_1(x, y) = x\left(x - 1\right)\left( \frac{2}{3}x - y + \frac{2}{3}\right) = \frac{2}{3}x + xy + \frac{2}{3}x^3-x^2y$.

 Another way to do this would to use the line $y = x$ and the ellipse, $x^2 + xy + y^2 = 1$. This would give

 $p_2(x, y) = \left(x - y\right)\left(x^2 + xy + y^2 - 1\right) = -x + y + x^3 - y^3$. (But 

 (Another way would be to use the hyperbola from Solution 1. Interesting that different curves both work.)

 At this point, we consider that $p_1$ and $p_2$ both must have $\left(\frac{a}{c}, \frac{b}{c}\right)$ as a zero. A quick graph of the 4 lines and the ellipse used to create $p_1$ and $p_2$ gives nine intersection points. Eight of them are the given ones, and the ninth is $\left(\frac{5}{19}, \frac{16}{19}\right)$. The last intersection point can be found by finding the intersection points of $y = \frac{2}{3}x + \frac{2}{3}$ and $x^2 + xy + y^2 = 1$.Finally, just add the values of $a$, $b$, and $c$ to get $5 + 16 + 19 = \boxed{040}$

                    \begin{figure}[H]
                        \centering
                        \includegraphics[width=0.5\textwidth]{Solution/AIME_I_2008_13_Solution_Figure_2.png}
                    \end{figure}
                    




\textbf{Solution 3 (Ansatz)}

We can plug in the values to obtain 

 
 \[p(0,0)=0\Longrightarrow a_0=0\]

 \[p(1,0)=0\Longrightarrow a_1+a_3+a_6=0\]

 \[p(0,1)=0\Longrightarrow a_2+a_5+a_9=0\]

 \[p(-1,0)=0\Longrightarrow a_1-a_3+a_6=0\]

 \[p(0,-1)=0\Longrightarrow a_2-a-5+a_9=0\]

 \[p(1,1)=0\Longrightarrow a_4+a_7+a_8=0\]

 \[p(1,-1)=0\Longrightarrow a_4+a_7-a_8=0\]

 \[p(2,2)=0\Longrightarrow2a_4=3a_6+3a_9\Leftrightarrow2a_4+3a_1+3a_2=0.\] 

 Now, this means that 

 
 \[p(x,y)=a_1x+a_2y+a_4xy-a_1x^3+a_7x^2y+a_8xy^2-a_2y^3.\] 

 After some simplifying, we obtain 

 
 \[p(x,y)=a_1(x-x^3)+a_2(y-y^3)+a_4xy(1-x).\] 

 Since $p(x,y)=0$, $3a_1+3a_2+2a_4=0,$ and we suspect that:

 
 \[x-x^3=y-y^3\] 

 and 

 
 \[\frac{x-x^3}{xy(1-x)}=\frac{3}{2}\Leftrightarrow\frac{1+x}{y}=\frac{3}{2} \Leftrightarrow\frac{2}{3}+\frac{2x}{3}=y\].

 Plugging this into the first equation, and factoring, and cancelling $(x+1)$, and simplifying, we get $19x^2 -43x+10=0$, so we find that $(x,y)=\left(\frac{5}{19},\frac{16}{19}\right)\Longrightarrow a+b+c=5+16+19=\boxed{40}.$



	\item (\textit{AIME I 2008}) Let $\overline{AB}$ be a diameter of circle $\omega$. Extend $\overline{AB}$ through $A$ to $C$. Point $T$ lies on $\omega$ so that line $CT$ is tangent to $\omega$. Point $P$ is the foot of the perpendicular from $A$ to line $CT$. Suppose $\overline{AB} = 18$, and let $m$ denote the maximum possible length of segment $BP$. Find $m^{2}$.


 
 \textbf{Answer: }432



\textbf{Solution 1}


                    \begin{figure}[H]
                        \centering
                        \includegraphics[width=0.5\textwidth]{Solution/AIME_I_2008_14_Solution_Figure_1.png}
                    \end{figure}
                    

 Let $x = OC$. Since $OT, AP \perp TC$, it follows easily that $\triangle APC \sim \triangle OTC$. Thus $\frac{AP}{OT} = \frac{CA}{CO} \Longrightarrow AP = \frac{9(x-9)}{x}$. By the \href{https://artofproblemsolving.com/wiki/index.php?title=Law\_of\_Cosines}{Law of Cosines} on $\triangle BAP$, 
 \begin{align*}BP^2 = AB^2 + AP^2 - 2 \cdot AB \cdot AP \cdot \cos \angle BAP \end{align*}
where $\cos \angle BAP = \cos (180 - \angle TOA) = - \frac{OT}{OC} = - \frac{9}{x}$, so: 
 \begin{align*}BP^2 &= 18^2 + \frac{9^2(x-9)^2}{x^2} + 2(18) \cdot \frac{9(x-9)}{x} \cdot \frac 9x = 405 + 729\left(\frac{2x - 27}{x^2}\right)\end{align*}
Let $k = \frac{2x-27}{x^2} \Longrightarrow kx^2 - 2x + 27 = 0$; this is a quadratic, and its \href{https://artofproblemsolving.com/wiki/index.php?title=Discriminant}{discriminant} must be nonnegative: $(-2)^2 - 4(k)(27) \ge 0 \Longleftrightarrow k \le \frac{1}{27}$. Thus,
 \[BP^2 \le 405 + 729 \cdot \frac{1}{27} = \boxed{432}\]Equality holds when $x = 27$.~Shen Kislay Kai

 



\textbf{Solution 1.1 (Calculus)}

Proceed as follows for Solution 1. 

 Once you approach the function $k=(2x-27)/x^2$, find the maximum value by setting $dk/dx=0$.

 Simplifying $k$ to take the derivative, we have $2/x-27/x^2$, so $dk/dx=-2/x^2+54/x^3$.  Setting $dk/dx=0$, we have $2/x^2=54/x^3$.  

 Solving, we obtain $x=27$ as the critical value.  Hence, $k$ has the maximum value of $(2*27-27)/27^2=1/27$.  Since $BP^2=405+729k$, the maximum value of $\overline {BP}$ occurs at $k=1/27$, so $BP^2$ has a maximum value of $405+729/27=\fbox{432}$.  

 Note: Please edit this solution if it feels inadequate.~Shen Kislay Kai

 



\textbf{Solution 2}


                    \begin{figure}[H]
                        \centering
                        \includegraphics[width=0.5\textwidth]{Solution/AIME_I_2008_14_Solution_Figure_2.png}
                    \end{figure}
                    

 From the diagram, we see that $BQ = OT + BO \sin\theta = 9 + 9\sin\theta = 9(1 + \sin\theta)$, and that $QP = BA\cos\theta = 18\cos\theta$.

 
 \begin{align*}BP^2 &= BQ^2 + QP^2 = 9^2(1 + \sin\theta)^2 + 18^2\cos^2\theta\\ &= 9^2[1 + 2\sin\theta + \sin^2\theta + 4(1 - \sin^2\theta)]\\ BP^2 &= 9^2[5 + 2\sin\theta - 3\sin^2\theta]\end{align*}


 This is a \href{https://artofproblemsolving.com/wiki/index.php?title=Quadratic\_equation}{quadratic equation}, maximized when $\sin\theta = \frac { - 2}{ - 6} = \frac {1}{3}$. Thus, $m^2 = 9^2[5 + \frac {2}{3} - \frac {1}{3}] = \boxed{432}$.

 



\textbf{Solution 3 (Calculus Bash)}


                    \begin{figure}[H]
                        \centering
                        \includegraphics[width=0.5\textwidth]{Solution/AIME_I_2008_14_Solution_Figure_3.png}
                    \end{figure}
                    

 (Diagram credit goes to Solution 2)

 We let $AC=x$. From similar triangles, we have that $PC=\frac{x\sqrt{x^2+18x}}{x+9}$ (Use Pythagorean on $\triangle\omega TC$ and then using $\triangle\omega CT\sim\triangle ACP$). Similarly, $TP=QT=\frac{9\sqrt{x^2+18x}}{x+9}$. Using the Pythagorean Theorem again and $\triangle CAP\sim\triangle CBQ$, $BQ=\sqrt{(x+18)^2-(\frac{(x+18)\sqrt{x^2+18x}}{x+9})^2}$. Using the Pythagorean Theorem $\bold{again}$, $BP=\sqrt{(x+18)^2-(\frac{(x+18)\sqrt{x^2+18x}}{x+9})^2+(\frac{18\sqrt{x^2+18x}}{x+9})^2}$. After a large bashful simplification, $BP=\sqrt{405+\frac{1458x-6561}{x^2+18x+81}}$. The fraction is equivalent to $729\frac{2x-9}{(x+9)^2}$. Taking the derivative of the fraction and solving for x, we get that $x=18$. Plugging $x=18$ back into the expression for $BP$ yields $\sqrt{432}$, so the answer is $(\sqrt{432})^2=\boxed{432}$.

 



\textbf{Solution 4}


                    \begin{figure}[H]
                        \centering
                        \includegraphics[width=0.5\textwidth]{Solution/AIME_I_2008_14_Solution_Figure_3.png}
                    \end{figure}
                    

 (Diagram credit goes to Solution 2)

 Let $AC=x$. The only constraint on $x$ is that it must be greater than $0$. Using similar triangles, we can deduce that $PA=\frac{9x}{x+9}$. Now, apply law of cosines on $\triangle PAB$. 
 \[BP^2=\left(\frac{9x^2}{x+9}\right)^2+18^2-2(18)\left(\frac{9x}{x+9}\right)\cos(\angle PAB).\] We can see that $\cos(\angle PAB)=\cos(180^{\circ}-\angle PAC)=\cos(\angle PAC -90^{\circ})=-\sin(\angle PCA)$. We can find $-\sin(\angle PCA)=-\frac{9}{x+9}$. Plugging this into our equation, we get: 
 \[BP^2=\left(\frac{9x^2}{x+9}\right)^2+18^2-2(18)\left(\frac{9x}{x+9}\right)\left(-\frac{9}{x+9}\right).\] Eventually, 
 \[BP^2 = 81\left(\frac{x^2+36x}{(x+9)^2}+4\right).\] We want to maximize $\frac{x^2+36x}{(x+9)^2}$. There are many ways to maximize this expression, discussed here: \href{https://artofproblemsolving.com/community/c4h2292700\_maximization}{https://artofproblemsolving.com/community/c4h2292700\_maximization}. The maximum result of that expression is $\frac{4}{3}$. Finally, evaluating $BP^2$ for this value $81\left(\frac{4}{3}+4\right) = \boxed{432}$. 

 



\textbf{Solution 5 (Clean)}

Let $h$ be the distance from $A$ to $CT$. Observe that $h$ takes any value from $0$ to $r$, where $r$ is the radius of the circle.

 Let $Q$ be the foot of the altitude from $B$ to $CT$. It is clear that $T$ is the midpoint of $PQ$, and so the length $OT$ is the average of $AP$ and $BQ$. It follows thus that $BQ = 2r - h$.

 We compute $PT = \sqrt{r^2 - (r - h)^2} = \sqrt{h(2r - h)},$and so $BP^2 = PQ^2 + BQ^2 = 4PT^2 + BQ^2 = 4h(2r - h) + (2r-h)^2 = (2r-h)(2r + 3h)$.This is $\frac{1}{3}(6r - 3h)(2r + 3h) \le \frac{1}{3} \cdot \left( \frac{8r}{2} \right)^2$. Equality is attained, so thus we extract the answer of $\frac{16 \cdot 9^2}{3} = 27 \cdot 16 = \boxed{432}.$




\textbf{Solution 6 (Calculus)}


                    \begin{figure}[H]
                        \centering
                        \includegraphics[width=0.5\textwidth]{Solution/AIME_I_2008_14_Solution_Figure_5.png}
                    \end{figure}
                    

 (Diagram Credits to Solution 2)

 We can see that $AP = \frac{9x}{x+9}$. From the Law of Cosines, we can now find the length of BP.

 $BP^2 = (\frac{9x}{x+9})^2 + 18^2 - 2(18)(\frac{9x}{x+9})(-\frac{9}{x+9})$
$BP^2 = 81(\frac{x^2}{(x+9)^2} + 4 + \frac{36x}{(x+9)^2})$
$BP^2 = 81(\frac{x^2+36x}{(x+9)^2}+4)$

 From here, we see that the value we need to maximize is $\frac{x^2+36x}{(x+9)^2}$. We shall now do some calculus. Let $f(x) = \frac{x^2+36x}{(x+9)^2}$. We need to find the first derivative. Using the quotient rule $(\frac{u}{v})' = \frac{u'v - uv'}{v^2}$, we can find the derivative which is:

 $f'(x) = \frac{(2x+36)(x+9)^2-(x^2+36x)(2)(x+9)}{((x+9)^2)^2}$

 This simplifies to

 $f'(x) = \frac{-18x+324}{(x+9)^3}$

 Now we find the critical point by setting $f'(x)$ to zero and solve for x.

 $-18x + 324 = 0$
$-18x = -324$
$x = 18$

 We now plug this value of x back into the function to find the max value of f(x)

 $f(18) = \frac{18^2+36(18)}{27^2}$
$f(18) = \frac{4}{3}$

 We now go back to the formula for $BP^2$ and plug in $\frac{4}{3}$ for that ugly fraction.

 $BP^2 = 81(\frac{4}{3}+4)$
$BP^2 = 81(\frac{16}{3})$
$BP^2 = 432$

 Thus, the value of $m^2$ (or $BP^2$) with the maximum value of $m$ is $\boxed{432}$.

 


	\item (\textit{AIME I 2008}) A square piece of paper has sides of length $100$. From each corner a wedge is cut in the following manner: at each corner, the two cuts for the wedge each start at distance $\sqrt {17}$ from the corner, and they meet on the diagonal at an angle of $60^\circ$ (see the figure below). The paper is then folded up along the lines joining the vertices of adjacent cuts. When the two edges of a cut meet, they are taped together. The result is a paper tray whose sides are not at right angles to the base. The height of the tray, that is, the perpendicular distance between the plane of the base and the plane formed by the upper edges, can be written in the form $\sqrt [n]{m}$, where $m$ and $n$ are positive integers, $m < 1000$, and $m$ is not divisible by the $n$th power of any prime. Find $m + n$.


 
                    \begin{figure}[H]
                        \centering
                        \includegraphics[width=0.5\textwidth]{Problem/AIME_I_2008_15_Problem_Figure_1.png}
                    \end{figure}
                    

 \textbf{Answer: }871



\textbf{Solution 1}

In the original picture, let $P$ be the corner, and $M$ and $N$ be the two points whose distance is $\sqrt{17}$ from $P$. Also, let $R$ be the point where the two cuts intersect.

 Using $\triangle{MNP}$ (a 45-45-90 triangle), $MN=MP\sqrt{2}\quad\Longrightarrow\quad MN=\sqrt{34}$. $\triangle{MNR}$ is \href{https://artofproblemsolving.com/wiki/index.php?title=Equilateral\_triangle}{equilateral}, so $MR = NR = \sqrt{34}$. (Alternatively, we could find this by the \href{https://artofproblemsolving.com/wiki/index.php?title=Law\_of\_Sines}{Law of Sines}.)

 The length of the perpendicular from $P$ to $MN$ in $\triangle{MNP}$ is $\frac{\sqrt{17}}{\sqrt{2}}$, and the length of the perpendicular from $R$ to $MN$ in $\triangle{MNR}$ is $\frac{\sqrt{51}}{\sqrt{2}}$. Adding those two lengths, $PR=\frac{\sqrt{17}+\sqrt{51}}{\sqrt{2}}$. (Alternatively, we could have used that $\sin 75^{\circ} = \sin (30+45) = \frac{\sqrt{6}+\sqrt{2}}{4}$.)

 Drop a \href{https://artofproblemsolving.com/wiki/index.php?title=Perpendicular}{perpendicular} from $R$ to the side of the square containing $M$ and let the intersection be $G$.

 
 \begin{align*}PG&=\frac{PR}{\sqrt{2}}=\frac{\sqrt{17}+\sqrt{51}}{2}\\ MG=PG-PM&=\frac{\sqrt{17}+\sqrt{51}}{2}-\sqrt{17}=\frac{\sqrt{51}-\sqrt{17}}{2}\end{align*}


                    \begin{figure}[H]
                        \centering
                        \includegraphics[width=0.5\textwidth]{Solution/AIME_I_2008_15_Solution_Figure_1.png}
                    \end{figure}
                    

 Let $A'B'C'D'$ be the smaller square base of the tray and let $ABCD$ be the larger square, such that $AA'$, etc, are edges. Let $F$ be the foot of the perpendicular from $A$ to plane $A'B'C'D'$.

 We know $AA'=MR=\sqrt{34}$ and $A'F=MG\sqrt{2}=\frac{\sqrt{51}-\sqrt{17}}{\sqrt{2}}$. Now, use the Pythagorean Theorem on triangle $AFA'$ to find $AF$:

 
 \begin{align*}\left(\frac{\sqrt{51}-\sqrt{17}}{\sqrt{2}}\right)^2+AF^2&=\left(\sqrt{34}\right)^2\\ \frac{51-34\sqrt{3}+17}{2}+AF^2&=34\\AF&=\sqrt{34-\frac{68-34\sqrt{3}}{2}}\\AF&=\sqrt{\frac{34\sqrt{3}}{2}}\\AF&=\sqrt[4]{867}\end{align*}


 The answer is $867 + 4 = \boxed{871}$.

 



\textbf{Solution 2}

In the final pyramid, let $ABCD$ be the smaller square and let $A'B'C'D'$ be the larger square such that $AA'$, etc. are edges.

 It is obvious from the diagram that $\angle A'AB = \angle A'AD = 105^\circ$.

 Let $AB$ and $AD$ be the positive $x$ and $y$ axes in a 3-d coordinate system such that $A'$ has a positive $z$ coordinate. Let $\alpha$ be the angle made with the positive $x$ \href{https://artofproblemsolving.com/wiki/index.php?title=Axis}{axis}. Define $\beta$ and $\gamma$ analogously.

 It is easy to see that if $P: = (x,y,z)$, then $x = AA'\cdot \cos\alpha$. Furthermore, this means that $\frac {x^2}{AA'^2} + \frac {y^2}{AA'^2} + \frac {z^2}{AA'^2} = \cos^2\alpha + \cos^2\beta + \cos^2\gamma = 1$.

 We have that $\alpha = \beta = 105^\circ$, so $\cos^2 105^\circ + \cos^2105^\circ + \cos^2\gamma = 1\implies \cos\gamma = \sqrt [4]{\frac {3}{4}}$.

 It is easy to see from the \href{https://artofproblemsolving.com/wiki/index.php?title=Law\_of\_Sines}{Law of Sines} that $\frac {AA'}{\sin 45^\circ} = \frac {\sqrt {17}}{\sin 30^\circ}\implies AA' = \sqrt {34}$.

 Now, $z = AA'\cdot \cos\gamma = \sqrt [4]{34^2\cdot \frac {3}{4}} = \sqrt [4]{867}$.

 It follows that the answer is $867 + 4 = \boxed{871}$.~Shen Kislay Kai

 



\textbf{Solution 3 (Vectors)}

Let $A$ be a corner of the square. Let $B$ and $C$ be the points where the cuts start from the edges, and $O$ be the point where the cuts meet. Let $X$ and $Y$ be points such that $OX=OY=1,$ $\overline{OX} \parallel \overline{AB}$ and $\overline{OY} \parallel \overline{AC},$ respectively. Let $P$ be the point such that the cuts $OB$ and $OC$ meet at $P$ when folded. 

 By the AAS congruence theorem, $\triangle{AOC} \cong \triangle{AOB}.$ Hence, it follows that $\triangle{OBC}$ is equilateral and $OB=OC=BC=OP.$ Since $\triangle{ABC}$ is an isosceles right triangle, we have $OP=OB=OC=BC=AB\cdot\sqrt{2}=\sqrt{17}\cdot\sqrt{2}=\sqrt{34}.$ Thus, $OP=\sqrt{34}.$

 Also, note that $\angle{POX}=\angle{BOX}=\angle{AOX}-\angle{AOB}=105^\circ.$

 Furthermore, let $O=(0,0,0)$ be the origin, $X=(1,0,0),$ and $Y=(0,1,0).$ Let $\overrightarrow{P}=\begin{pmatrix} p_1 \\ p_2 \\ p_3 \end{pmatrix}$ denote the vector at $P,$ $\overrightarrow{X}=\begin{pmatrix} 1 \\ 0 \\ 0 \end{pmatrix}$ denote the vector at $X,$ and $\overrightarrow{Y}=\begin{pmatrix} 0 \\ 1 \\ 0 \end{pmatrix}$ denote the vector at $Y.$ It suffices to solve for $p_3.$

 By the properties of the dot product for vectors, 
 \[\overrightarrow{P}\cdot\overrightarrow{X}=\overrightarrow{P}\cdot\overrightarrow{Y}=\|\overrightarrow{P}\|\|\overrightarrow{X}\|\cos\angle{POX}=\|\overrightarrow{P}\|\|\overrightarrow{Y}\|\cos\angle{POY}.\] Since $\angle{POX}=\angle{POY}=105^\circ,$ $\|\overrightarrow{P}\|=OP=\sqrt{34},$ and $\|\overrightarrow{X}\|=\|\overrightarrow{Y}\|=1,$ the above equations simplify to
 \[\begin{pmatrix} p_1 \\ p_2 \\ p_3 \end{pmatrix}\cdot\begin{pmatrix} 1 \\ 0 \\ 0 \end{pmatrix}=\begin{pmatrix} p_1 \\ p_2 \\ p_3 \end{pmatrix}\cdot\begin{pmatrix} 0 \\ 1 \\ 0 \end{pmatrix}=\sqrt{34}\cos 105^\circ.\]This leads to $p_1=p_2=\sqrt{34}\cos 105^\circ.$ Also, since $\|\overrightarrow{P}\|=\sqrt{34},$ we have \begin{align*}p_3^2&=34-(p_1^2+p_2^2)\\&=34-2\cdot(\sqrt{34}\cos 105^\circ)^2\\&=34(1-2\cos^2 105^\circ)\\&=-34\cos 210^\circ\\&=34\cos 30^\circ\\&=34\frac{\sqrt{3}}{2}\\&=17\sqrt{3}\\&=\sqrt{867}.\\\end{align*}


 Thus, $p_3=\sqrt[4]{867}$ and $m+n=\boxed{871}.$



\end{enumerate}